\section*{Erd\H{o}s Problem \#276}

\subsection*{1) ROUND-2 OBJECTIVE}
We pursue path \textbf{(C): obstruction/correction}. Round~1 already separated the weak Graham--Vsemirnov statement from the stronger Erd\H{o}s problem, but it left one conceptual gap: the informal phrase
\begin{quote}
``without an underlying system of covering congruences responsible''
\end{quote}
was not yet formalized exactly.

The goal of Round~2 is to close that gap. We prove that for the recurrence
\[
a_{n+2}=a_{n+1}+a_n,
\]
the natural precise meaning of ``a finite covering system is responsible'' is \emph{equivalent} to the literal strong condition
\[
\forall m>1\ \exists n\ge 0:\ \gcd(m,a_n)=1.
\]
So the intended conjecture is \emph{not} a different one: after the minimal natural formalization, it is the same conjecture.

\subsection*{2) Round-1 FOUNDATION USED}
We use the following Round~1 results without re-proving them.

\begin{itemize}
\item[(F1)] The weak statement is solved: there exist relatively prime positive integers $a_0,a_1$ such that every term of the corresponding Lucas sequence is composite (Graham; record-small known example due to Vsemirnov).
\item[(F2)] Any covering-congruence construction of Graham type cannot solve the strong Erd\H{o}s problem, because the product of the finitely many covering primes has a nontrivial gcd with every term.
\item[(F3)] For Fibonacci numbers $F_n$, one has the addition formula
\[
F_{u+v}=F_uF_{v+1}+F_{u-1}F_v
\]
and the divisibility property
\[
F_m\mid F_{tm}\qquad (m,t\ge 1).
\]
\item[(F4)] For a Fibonacci-like sequence, the gcd of consecutive terms is constant:
\[
\gcd(a_n,a_{n+1})=\gcd(a_0,a_1)\qquad (n\ge 0).
\]
In particular, any sequence satisfying the strong Erd\H{o}s condition automatically has $\gcd(a_0,a_1)=1$.
\end{itemize}

\subsection*{3) NEW INSIGHT / TOOL (ROUND-2)}
The new structural fact is the following.

\medskip
\noindent\textbf{Single-class zero theorem modulo a prime.}
Fix a prime $p$ and a Fibonacci-like sequence $(a_n)_{n\ge 0}$. Then the set
\[
Z_p:=\{n\ge 0: p\mid a_n\}
\]
is of one of the following three types:
\begin{enumerate}
\item[(i)] $Z_p=\varnothing$;
\item[(ii)] $Z_p=\{0,1,2,\dots\}$;
\item[(iii)] $Z_p$ is a \emph{single} arithmetic progression
\[
Z_p=\{n\ge 0: n\equiv r_p\pmod{z(p)}\},
\]
where $z(p)$ is the least positive integer with $p\mid F_{z(p)}$.
\end{enumerate}
For sequences relevant to the strong Erd\H{o}s problem, case (ii) cannot occur because then $p$ would divide $a_0$ and $a_1$, contradicting $\gcd(a_0,a_1)=1$.

Consequently, failure of the strong condition
\[
\forall m>1\ \exists n:\ \gcd(m,a_n)=1
\]
is \emph{equivalent} to the existence of a finite family of primes whose corresponding single congruence classes cover all indices. This gives the exact mathematical meaning of ``an underlying covering system of congruences responsible.''

\subsection*{4) ATTACK PLAN (ROUND-2)}
Round~1 already showed that the strong problem was not settled by Graham/Vsemirnov-type constructions. The exact remaining gap was:
\begin{quote}
Does the literal Erd\H{o}s condition merely \emph{suggest} the absence of a covering system, or is it actually \emph{equivalent} to the absence of a finite prime-based congruence covering?
\end{quote}
To close this gap, we prove four claims.
\begin{enumerate}
\item[(A)] For every prime $p$, the Fibonacci numbers have a well-defined rank of apparition $z(p)$, and
\[
p\mid F_n \iff z(p)\mid n.
\]
\item[(B)] For every Fibonacci-like sequence and every prime $p$, the set of indices where $p$ divides the sequence is either empty, all of $\mathbf{N}_0$, or a single residue class modulo $z(p)$.
\item[(C)] Therefore a finite common-factor obstruction
\[
\exists m>1\ \forall n:\ \gcd(m,a_n)>1
\]
is exactly the same as a finite prime-covering congruence system.
\item[(D)] For any \emph{fixed} finite set of primes, deciding whether that set covers the sequence reduces to a finite computation.
\end{enumerate}
This is strict progress beyond Round~1 because it turns the informal ``covering congruences responsible'' phrase into a rigorous equivalent condition.

\subsection*{5) WORK (ROUND-2)}

Let $F_0=0$, $F_1=1$, and $F_{n+2}=F_{n+1}+F_n$ be the Fibonacci numbers.

\medskip
\noindent\textbf{Claim A.} For every prime $p$, there exists a least positive integer $z(p)$ such that $p\mid F_{z(p)}$. Moreover,
\[
p\mid F_n \iff z(p)\mid n \qquad (n\ge 0).
\]

\medskip
\noindent\emph{Proof.}
Let
\[
A=\left(\begin{array}{cc}0&1\\ 1&1\end{array}\right).
\]
Modulo $p$, the matrix $A$ lies in the finite group $\mathrm{GL}_2(\mathbf{F}_p)$ because $\det A=-1\not\equiv 0\pmod p$. Hence $A$ has finite order $T_p>0$. Since
\[
A^n\left(\begin{array}{c}0\\ 1\end{array}\right)=
\left(\begin{array}{c}F_n\\ F_{n+1}\end{array}\right),
\]
we get
\[
A^{T_p}\left(\begin{array}{c}0\\ 1\end{array}\right)=
\left(\begin{array}{c}0\\ 1\end{array}\right),
\]
so $p\mid F_{T_p}$. Therefore the least positive such index $z(p)$ exists.

If $z(p)\mid n$, then by Foundation~(F3) we have $F_{z(p)}\mid F_n$, hence $p\mid F_n$.

Conversely, suppose $p\mid F_n$. If $n<z(p)$, then by minimality of $z(p)$ we must have $n=0$, so certainly $z(p)\mid n$. Now assume $n\ge z(p)$ and write
\[
n=qz(p)+r,\qquad q\ge 1,\qquad 0\le r<z(p).
\]
By the addition formula from Foundation~(F3),
\[
F_n=F_{qz(p)+r}=F_{qz(p)}F_{r+1}+F_{qz(p)-1}F_r.
\]
Since $z(p)\mid qz(p)$, the first term is $0\pmod p$, so
\[
F_n\equiv F_{qz(p)-1}F_r\pmod p.
\]
Now use Cassini's identity:
\[
F_{qz(p)+1}F_{qz(p)-1}-F_{qz(p)}^2=(-1)^{qz(p)}.
\]
Because $p\mid F_{qz(p)}$, this gives
\[
F_{qz(p)+1}F_{qz(p)-1}\equiv (-1)^{qz(p)}\not\equiv 0\pmod p,
\]
so $F_{qz(p)-1}\not\equiv 0\pmod p$. Therefore $p\mid F_n$ implies $p\mid F_r$. By the minimality of $z(p)$, this forces $r=0$. Hence $z(p)\mid n$.

This proves Claim~A.
\hfill$\square$

\medskip
\noindent\textbf{Claim B.} Let $(a_n)_{n\ge 0}$ satisfy $a_{n+2}=a_{n+1}+a_n$, and fix a prime $p$. Then the set
\[
Z_p=\{n\ge 0: p\mid a_n\}
\]
is either empty, all of $\mathbf{N}_0$, or a single residue class modulo $z(p)$.

\medskip
\noindent\emph{Proof.}
First note that if $p\mid a_r$ and $p\mid a_{r+1}$ for some $r$, then the recurrence gives $p\mid a_{r+2}$, and the reversible form
\[
a_{r-1}=a_{r+1}-a_r
\]
gives $p\mid a_{r-1}$. Inducting forward and backward, $p$ divides every term. This is case~(ii).

Now suppose case~(ii) does \emph{not} hold, but $Z_p\neq \varnothing$. Pick $r\ge 0$ with $p\mid a_r$. Then necessarily $p\nmid a_{r+1}$, for otherwise $p$ would divide all terms.

Work modulo $p$, and extend the recurrence uniquely to all integers by the rule
\[
a_{n-1}=a_{n+1}-a_n.
\]
Also extend the Fibonacci sequence to all integers by
\[
F_{-t}=(-1)^{t+1}F_t\qquad (t\ge 1),
\]
so that the recurrence still holds for every integer index.

Let $b$ be the residue class of $a_{r+1}$ modulo $p$. Consider the bi-infinite sequence
\[
u_n:=bF_{n-r}\qquad (n\in \mathbf{Z}).
\]
It satisfies the same recurrence as $(a_n)$ modulo $p$, and
\[
u_r=bF_0=0\equiv a_r\pmod p,
\]
\[
u_{r+1}=bF_1=b\equiv a_{r+1}\pmod p.
\]
Since a second-order linear recurrence is uniquely determined by two consecutive values, we have
\[
a_n\equiv u_n\equiv bF_{n-r}\pmod p
\qquad (n\in \mathbf{Z}).
\]
Because $b\not\equiv 0\pmod p$, this yields
\[
p\mid a_n \iff p\mid F_{n-r}.
\]
By Claim~A,
\[
p\mid F_{n-r} \iff z(p)\mid (n-r).
\]
Therefore
\[
Z_p=\{n\ge 0: n\equiv r\pmod{z(p)}\}.
\]
So in the nontrivial case, $Z_p$ is a single residue class modulo $z(p)$.

This proves Claim~B.
\hfill$\square$

\medskip
\noindent\textbf{Definition.} Let $(a_n)$ be a Fibonacci-like sequence with $\gcd(a_0,a_1)=1$. A \emph{finite prime-covering congruence system} for $(a_n)$ means a finite set of primes $p_1,\dots,p_t$ and residues $r_i$ modulo $z(p_i)$ such that
\[
p_i\mid a_n \iff n\equiv r_i\pmod{z(p_i)}
\qquad (1\le i\le t),
\]
and the residue classes
\[
n\equiv r_i\pmod{z(p_i)}
\]
cover all integers $n\ge 0$.

\medskip
\noindent\textbf{Claim C.} Assume $\gcd(a_0,a_1)=1$. Then the following are equivalent.
\begin{enumerate}
\item[(1)] There exists an integer $m>1$ such that
\[
\gcd(m,a_n)>1\qquad \text{for all } n\ge 0.
\]
\item[(2)] The sequence $(a_n)$ admits a finite prime-covering congruence system.
\end{enumerate}

\medskip
\noindent\emph{Proof.}
Assume (1). Let $S$ be the finite set of prime divisors of $m$. Because $\gcd(a_0,a_1)=1$, no prime can divide every term of the sequence. Hence for each $p\in S$, Claim~B applies in the nontrivial case: either $p$ divides no term, or there is a unique residue class
\[
n\equiv r_p\pmod{z(p)}
\]
on which $p$ divides $a_n$. Since $\gcd(m,a_n)>1$ for every $n$, for each $n$ some prime $p\in S$ divides $a_n$. Thus the finitely many residue classes coming from those primes cover all $n\ge 0$. This is exactly (2).

Conversely, assume (2), and set
\[
m:=\prod_{i=1}^t p_i.
\]
Given any $n\ge 0$, the covering property yields an index $i$ with
\[
n\equiv r_i\pmod{z(p_i)}.
\]
By definition of the prime-covering congruence system, $p_i\mid a_n$. Therefore
\[
\gcd(m,a_n)\ge p_i>1.
\]
So $\gcd(m,a_n)>1$ for every $n\ge 0$, proving (1).

This proves Claim~C.
\hfill$\square$

\medskip
\noindent\textbf{Claim D.} Let $S$ be a fixed finite set of primes. Then deciding whether
\[
\forall n\ge 0\ \exists p\in S:\ p\mid a_n
\]
holds is a finite computation.

\medskip
\noindent\emph{Proof.}
For each $p\in S$, Claim~B shows that the set of indices $n$ with $p\mid a_n$ is either empty or a single class modulo $z(p)$ (under the nontrivial coprime assumption relevant to the strong problem). Let
\[
L:=\mathrm{lcm}\{z(p): p\in S\}.
\]
Then the condition
\[
\forall n\ge 0\ \exists p\in S:\ p\mid a_n
\]
is equivalent to the statement that every residue class modulo $L$ lies in at least one of the lifted congruence classes coming from the primes in $S$. This requires checking only the finitely many residues $0,1,\dots,L-1$.

This proves Claim~D.
\hfill$\square$

\medskip
\noindent\textbf{Consequence.} The literal Erd\H{o}s condition
\[
\forall m>1\ \exists n\ge 0:\ \gcd(m,a_n)=1
\]
is exactly equivalent to saying that \emph{no finite prime-covering congruence system exists}. Hence the phrase
\begin{quote}
``without an underlying system of covering congruences responsible''
\end{quote}
has a canonical precise formalization for this problem, and that formalization is the strong Erd\H{o}s statement itself.

\subsection*{6) ADVERSARIAL VERIFICATION}
We check the new work against the main possible failure modes.

\begin{itemize}
\item \textbf{Could $Z_p$ be a union of two residue classes?} No. Claim~B shows that once a single zero $a_r\equiv 0\pmod p$ occurs and $p$ does not divide all terms, the entire sequence modulo $p$ is forced to be a scalar multiple of $F_{n-r}$, so the zero set is exactly one class modulo $z(p)$.

\item \textbf{Did the proof of Claim~A assume a theorem not justified here?} No. Existence of $z(p)$ came from the finite order of the companion matrix in $\mathrm{GL}_2(\mathbf{F}_p)$; the divisibility criterion $p\mid F_n\iff z(p)\mid n$ was proved directly from the Round~1 addition formula, the Round~1 divisibility lemma, and Cassini's identity.

\item \textbf{What about the case that a prime divides all terms?} Claim~B isolates that as a separate case. In the actual Erd\H{o}s problem, any admissible sequence must satisfy $\gcd(a_0,a_1)=1$, so this case cannot occur.

\item \textbf{Did Claim~C accidentally use the weak statement instead of the strong one?} No. Claim~C is purely structural: it does not assume all terms are composite. It only compares a fixed finite gcd-obstruction with a finite prime-covering congruence system.

\item \textbf{Does Claim~D solve the whole problem by finite computation?} No. It only says that for a \emph{prescribed finite set of primes} $S$, checking whether $S$ covers the sequence is finite. The open difficulty is that there is no known a priori bound on which primes could appear in such a set.
\end{itemize}

No contradiction with the Round~1 foundation appears.

\subsection*{7) FINAL}
\textbf{UNRESOLVED (BUT STRICTLY ADVANCED).}

Round~2 proves a new equivalence that was not available in Round~1:

\begin{quote}
For Fibonacci-like sequences, the literal strong Erd\H{o}s condition
\[
\forall m>1\ \exists n\ge 0:\ \gcd(m,a_n)=1
\]
is \emph{equivalent} to the absence of any finite prime-covering congruence system.
\end{quote}

Therefore the intended actual conjecture is \emph{not} different from the literal one. After the minimal natural formalization of ``an underlying covering system is responsible'', they are the same statement.

This does \emph{not} solve Problem~\#276, but it sharpens the remaining obstruction to one precise gap:
\begin{quote}
Construct an all-composite Fibonacci-like sequence for which no finite family of primes gives a covering by the single congruence classes $n\equiv r_p\pmod{z(p)}$.
\end{quote}

Equivalently, one must prove that for every finite set of primes $S$, there exists an index $n$ such that no prime in $S$ divides $a_n$.

\subsection*{8) COMPLETION ESTIMATE}
COMPLETION: 80\%.

\subsection*{9) REFERENCES}
\begin{itemize}
\item R. L. Graham, \emph{A Fibonacci-Like Sequence of Composite Numbers}, Mathematics Magazine 37 (1964), 322--324.
\item M. Vsemirnov, \emph{A New Fibonacci-Like Sequence of Composite Numbers}, Journal of Integer Sequences 7 (2004), Article 04.3.7.
\item D. Ismailescu and J. Son, \emph{A New Kind of Fibonacci-Like Sequence of Composite Numbers}, Journal of Integer Sequences 17 (2014), Article 14.8.2.
\item T. F. Bloom, \emph{Erd\H{o}s Problem \#276}, erdosproblems.com/276, accessed 2026-03-07.
\item T. F. Bloom, \emph{Erd\H{o}s Problem \#276: discussion thread}, erdosproblems.com/forum/thread/276, accessed 2026-03-07.
\end{itemize}
